\chapter{Axioms of the Real Numbers and Supremum}

A program searches for the length of the diagonal of a square with side one. It keeps
two numbers, $a=1$ and $b=2$, with $a^2-2<0<b^2-2$. At every step it computes the
midpoint $m=(a+b)/2$, evaluates $m^2-2$, and then replaces $a$ or $b$ by $m$ so that
the signs at the two ends stay different. After 40 steps the width of the interval is
below $10^{-12}$, after 100 steps below $10^{-30}$. But the program never produces an
answer: no step yields $x^2-2=0$ exactly, and the two endpoints keep squeezing
something that is not a rational number.

That failure is not a mistake in the algorithm. It shows that the rational numbers
have holes, and that the set $\{q\in\mathbb{Q}: q^2<2\}$ has no least upper bound
inside $\mathbb{Q}$.

Chapter 1 taught us to write proofs. Chapter 2 gave us the language of sets and
functions, and used $|x|$ and $\sup$ as computing tools without asking where they come
from. What is still missing is the object that makes those tools legitimate. This
chapter builds it, not out of decimals, but out of three groups of statements accepted
without proof: the field laws, the order laws, and completeness. Everything else,
including the existence of $\sqrt{2}$, is derived from there.

\begin{untukanak}
Imagine a ruler that has only fraction marks: one half, one quarter, one eighth, and
ever closer together. There are no other numbers on that ruler. Now measure the
diagonal of a square whose side is one hand-span. The end of the diagonal falls
between two marks, and it always does, no matter how close the marks are. We can get
as near as we like, but never exactly.

The question is: does that length really exist? The ruler clearly has no answer. Yet
the diagonal, if you carry it around or lay it on the floor, has a definite length.
That length does not disappear merely because our measuring tool does not contain it.

The same story happens with numbers. If we have only fractional numbers, there is a
point on the line that has no name. Mathematicians close that hole with a single
promise: every collection of numbers that is bounded above has a lowest roof. In this
chapter we will learn why that promise is needed, what its consequences are, and how
it gives a name to the diagonal length above, namely root two.
\end{untukanak}

\section{Field axioms}
\tingkat{basic}

We do not define the real numbers as endless decimals: that approach makes addition
complicated, and makes the definition of a least upper bound circular, since
``the supremum of all decimal cuts'' already uses the numbers being built. It is
cleaner to state the properties first and then derive the rest.

The first layer is algebraic. We assume two operations that obey the rules we have
used since school; the following definition only gives those rules a name.

\begin{definisi}[Field]\index{field}
A set $F$ is called a \emph{field} if addition and multiplication are defined on its
members so that: the result of both operations always lies in $F$; both operations are
commutative and associative; there is an element $0$ with $0+x=x$ and an element
$1\neq 0$ with $1\cdot x=x$; every $x$ has a negative $-x$ with $x+(-x)=0$; every
$x\neq 0$ has a reciprocal $1/x$ with $x\cdot(1/x)=1$; and $x(y+z)=xy+xz$ holds.
\end{definisi}

Those axioms mention neither order, nor distance, nor completeness. Other sets
therefore satisfy them too. The rational numbers with ordinary addition and
multiplication form a field. A reasonable question: what is already contained inside
those rules?

Three algebraic statements we usually memorize turn out to be consequences, not
axioms; the one most often used is multiplication by zero, which destroys everything.

\begin{teorema}
In an arbitrary field, for all $x,y$: (a) $0\cdot x=0$;
(b) $(-x)y=-(xy)$; (c) $(-x)(-y)=xy$; (d) if $xy=0$ then $x=0$ or $y=0$.
\end{teorema}

\begin{proof}
(a) From $0+0=0$ and the distributive law, $0\cdot x=(0+0)x=0\cdot x+0\cdot x$.
Add $-(0\cdot x)$ to both sides. The left side becomes $0$, the right side becomes
$0\cdot x$. Hence $0=0\cdot x$.

(b) Compute $(-x)y+xy=(-x+x)y=0\cdot y=0$, using (a). Hence $(-x)y$ is the negative
of $xy$, and that negative is unique.

(c) Apply (b) twice: $(-x)(-y)=-\bigl(x(-y)\bigr)=-\bigl(-(xy)\bigr)=xy$.
The last step uses $--a=a$, which is also a consequence of the uniqueness of the
negative.

(d) Suppose $x\neq 0$ and $xy=0$. Multiply by $1/x$: $(1/x)(xy)=0$, so that
$y=0$. Hence one of the factors is zero.
\end{proof}

Not one of those four statements is assumed; all of them descend from the distributive
law and the existence of negatives. This pattern is used throughout the book: state as
few rules as possible, then harvest the consequences.

\kamus{field}{lapangan}{A set with two arithmetic operations whose rules stay fixed, like
a toolbox whose lid always shuts tight: the result of a computation never leaves the
box.}

Chapter 2 used $\mathbb{Q}$ as a countable set. Here it appears in a different role:
as an ordered field. Those two properties are not yet enough, and that is the content
of the next section.

\section{Order axioms}
\tingkat{basic}

A field alone cannot compare two numbers. We need a relation ``less than'' that
respects addition and multiplication, because without it inequalities cannot be
manipulated. The order axioms say three things: any two numbers can be compared, the
comparison is transitive, and adding or multiplying by a positive number does not
break it.

\begin{definisi}[Ordered field]\index{ordered field}
A field $F$ is called \emph{ordered} if there is a relation $<$ on $F$ such that for
all $a,b,c$: exactly one of $a<b$, $a=b$, $b<a$ holds; if $a<b$ and $b<c$ then
$a<c$; if $a<b$ then $a+c<b+c$; and if $a<b$ and $c>0$ then $ac<bc$.
\end{definisi}

Almost every sign rule we memorize is a consequence of that definition, not an
addition to it. The following proof derives five of them at once; the pattern is the
same: add something to both sides, or multiply by a number whose sign is known.

\begin{proposisi}
In every ordered field: (a) if $x>0$ then $-x<0$, and conversely;
(b) if $x>0$ and $y<z$ then $xy<xz$; (c) if $x<0$ and $y<z$ then $xy>xz$;
(d) if $x\neq 0$ then $x^2>0$, in particular $1>0$; (e) if $0<x<y$ then
$0<1/y<1/x$.
\end{proposisi}

\begin{proof}
(a) If $x>0$, add $-x$ to both sides: $0>-x$. Conversely, if $x<0$ then
$0-x>0$, that is, $-x>0$.

(b) Since $z>y$, we have $z-y>0$. Multiply by $x>0$: $x(z-y)>0$. Hence
$xz=x(z-y)+xy>xy$.

(c) From (a), $-x>0$. Apply (b) to $-x$ and $y<z$: $(-x)y<(-x)z$. By the previous
Theorem, the left side is $-(xy)$ and the right side is $-(xz)$. Hence
$-(xy)<-(xz)$, and by (a) again $xy>xz$.

(d) If $x>0$, the order axiom applied to $y=x>0$ gives $x^2>0$. If $x<0$ then
$-x>0$, so that $(-x)^2>0$; yet $(-x)^2=x^2$. For $x=0$ clearly $x^2=0$. Since
$1=1^2$, we obtain $1>0$.

(e) Since $y>0$ and $1/y\neq 0$, from (d) $1/y>0$; the same holds for $1/x$.
Multiply $x<y$ by the positive number $(1/x)(1/y)$: the result is $1/y<1/x$.
\end{proof}

Two consequences that are used often: $a<b$ is equivalent to $b-a>0$, and $a<b$ is
equivalent to $-b<-a$; the second explains why multiplying by a negative number
reverses the direction, because multiplying by $-1$ reflects the number line about
zero. Another consequence: if $0<a<b$ then $a^2<b^2$, simply by multiplying $a<b$ by
$a>0$ and then by $b>0$ and adding the two. Hence comparing squares is equivalent to
comparing the original positive numbers; every time we write $q^2<2$ for a rational
$q$, we are comparing $q$ with something that need not be rational.

\kamus{ordered field}{lapangan terurut}{A field whose members can be arranged like the
steps of a staircase, and whose arithmetic rules do not disturb that arrangement.}

Note what we do not have yet. We do not know whether there is a positive number smaller
than every $1/n$, and we do not know whether $\{q\in\mathbb{Q}: q^2<2\}$ has a roof.
That is not a question of order but a question of completeness.

\section{Absolute value and the triangle inequality}
\tingkat{intermediate}

In Chapter 2 we wrote $|x-y|$ as a distance, and used it to compare errors, without
proving that the symbol behaves like a distance. Now its properties are derived from
the order. Absolute value discards the sign and leaves the distance to the point zero;
the following definition writes two branches, and it is the second branch that is
usually used.

\begin{definisi}[Absolute value]\index{absolute value}
For $x\in\mathbb{R}$, $|x|=x$ if $x\ge 0$ and $|x|=-x$ if $x<0$. Equivalently,
$|x|=\max\{x,-x\}$, and for real numbers $|x|=\sqrt{x^2}$ with the root discussed in
the next section.
\end{definisi}

From that definition it follows directly: $|x|\ge 0$, $|{-x}|=|x|$, $|x|^2=x^2$, and
$|xy|=|x|\,|y|$. The last is proved by checking the four sign combinations, or by
noting that both sides are nonnegative and their squares agree.

Inequalities containing an absolute value are easiest to handle after being converted
into two ordinary inequalities. The following lemma is the reason that conversion is
legitimate.

\begin{lemma}
For $c\ge 0$, $|x|\le c$ if and only if $-c\le x\le c$.
\end{lemma}

\begin{proof}
If $|x|\le c$: for $x\ge 0$ we have $x=|x|\le c$ and $-c\le 0\le x$; for
$x<0$ we have $-x=|x|\le c$, so that $x\ge -c$, and $x<0\le c$. Hence
$-c\le x\le c$.

Conversely, if $-c\le x\le c$: for $x\ge 0$, $|x|=x\le c$; for $x<0$,
$|x|=-x\le c$ because $x\ge -c$. In both cases $|x|\le c$.
\end{proof}

The triangle inequality is the only tool we need to control the difference of two
numbers. Its name comes from a picture: the length of one side of a triangle does not
exceed the sum of the lengths of the other two sides.

\begin{teorema}[Triangle inequality]\index{triangle inequality}
For all $x,y\in\mathbb{R}$, $|x+y|\le |x|+|y|$. Moreover
$\bigl||x|-|y|\bigr|\le |x-y|$.
\end{teorema}

\begin{proof}
From the definition, $-|x|\le x\le |x|$ and $-|y|\le y\le |y|$. Add the two pairs:
$-(|x|+|y|)\le x+y\le |x|+|y|$. By the previous lemma with $c=|x|+|y|$, we obtain
$|x+y|\le |x|+|y|$.

For the second part, write $x=(x-y)+y$ and apply the first part:
$|x|\le |x-y|+|y|$, so that $|x|-|y|\le |x-y|$. Swapping the roles of $x$ and $y$
gives $|y|-|x|\le |x-y|$. Since $c\le d$ and $-c\le d$ is equivalent to $|c|\le d$,
we have $\bigl||x|-|y|\bigr|\le |x-y|$.
\end{proof}

As a check, take $x=3$ and $y=-5$: $|x+y|=2$ while $|x|+|y|=8$, so the inequality is
slack. For $x=3$ and $y=5$ both sides equal $8$ and the inequality becomes an equation;
that happens exactly when the signs agree or one of them is zero.

Its practical use: bounds on distance become bounds on functions. If
$|x-\sqrt2|<\delta$, then
\[
|x^2-2|=|x-\sqrt2|\,|x+\sqrt2|<\delta\,(2\sqrt2+\delta),
\]
because $|x+\sqrt2|\le|x-\sqrt2|+2\sqrt2$. For $\delta=10^{-3}$ that bound equals
$0.002829427$, and the true value at $x=\sqrt2+10^{-3}$ is also $0.002829427$ since
$x^2-2=2\sqrt2\,\delta+\delta^2$ exactly: the bound is tight. For $\delta=10^{-6}$ the
same bound shrinks to $2.828428\times10^{-6}$.

\kamus{absolute value}{nilai mutlak}{The distance of a point from zero on the number line,
regardless of which side of zero the point lies on.}

The triangle inequality uses only the order axioms, not completeness, so it also holds
in $\mathbb{Q}$ and cannot distinguish the two. The distinguishing tool appears in the
next section.

\section{Completeness: the least upper bound}
\tingkat{intermediate}

Order gives comparison, but it does not give a roof. The set
$A=\{1-1/n: n\in\mathbb{N}\}$ contains numbers that come ever closer to one without
ever equaling one: one is not a member, yet it is the tightest bound. The set
$B=\{q\in\mathbb{Q}: q^2<2\}$ is bounded, yet it has no tight bound in
$\mathbb{Q}$. Both show that a slack bound and a tightest bound need different names.

\begin{untukanak}
Imagine we must put up a tent roof for a group of children. A roof that is too low:
some head sticks out. A roof that is too high: the tent wastes material and wobbles.
What we are looking for is the lowest roof that still lets all the children fit. Call
it the right roof.

Now something strange happens. In that group there is no tallest child; there is always
someone a little taller. Yet the right roof still exists, because their heights are
bounded: nobody is as tall as a tree. The right roof is determined by the whole group,
not by a single child.

In this chapter the right roof has a name: the least upper bound. The property that
guarantees a right roof always exists is called completeness. The fractional numbers
are incomplete: there are groups of bounded height whose right roof is not among the
fractional numbers. A question you can take home: if there is no tallest child, how do
you prove that the right roof really is the lowest one?
\end{untukanak}

An upper bound is a ceiling that contains every member; the least upper bound is the
lowest ceiling among them. The following definition names both, together with their
mirror image below.

\begin{definisi}[Upper bound and supremum]\index{upper bound}\index{supremum}
Let $A\subseteq\mathbb{R}$ and $A\neq\varnothing$. A number $u$ is called an
\emph{upper bound} of $A$ if $a\le u$ for every $a\in A$. If $A$ has at least one upper
bound, $A$ is called bounded above. A number $s$ is called the \emph{supremum} of $A$,
written $s=\sup A$, if $s$ is an upper bound of $A$ and $s\le u$ for every upper bound
$u$. Dually, $t$ is called the \emph{infimum} of $A$, written $t=\inf A$, if $t$ is a
lower bound and $t\ge v$ for every lower bound $v$.
\end{definisi}

The two conditions in the definition of supremum are often confused: the first says
that $s$ is still an upper bound, the second that no upper bound is smaller. For
$A=\{1-1/n\}$, the number $1$ satisfies both while $2$ satisfies only the first; that
set also has no largest member, so ``maximum'' and ``supremum'' are not synonyms.

\begin{bacakeras}
Two voices. The first voice asks, the second answers. Read them aloud alternately with
a friend.

``What is that fence?'' ``The fence is an upper bound: every member of the collection
lies below it.'' ``How many fences are there?'' ``Many. Every higher fence is also a
fence.'' ``Then which one do we use?'' ``The lowest one. Its name is the least upper
bound.'' ``What if that lowest fence happens to be occupied by a member of the
collection?'' ``Then it is also a maximum. If it is not occupied, it is still the least
upper bound. Two names, two jobs.''

Repeat three times, faster each time: a tall fence is not the answer, the lowest fence
is the answer.

The line the reader must repeat, once, in a clear voice: an upper bound is a fence; the
least upper bound is the lowest fence that still holds all the members in.
\end{bacakeras}

Figure~\ref{fig:sup-garis} places both ideas on a single line. The members of
$A=\{1-1/n\}$ crowd toward one direction without ever reaching the end; every number to
the right of one is an upper bound, and the leftmost among them is the supremum.

\begin{figure}[htbp]
\centering
\begin{tikzpicture}[x=1.5cm,y=1cm,font=\small]
  \draw[-{Latex[length=2mm]},ink,thick] (-0.35,0) -- (5.5,0);
  \foreach \t in {0,1,2,3,4,5}{
    \draw[ink] (\t,0.07) -- (\t,-0.07);
    \node[below=2pt,ink] at (\t,0) {$\t$};
  }
  \foreach \n in {1,...,6}{
    \fill[accent] ({1-1/\n},1.0) circle (1.7pt);
  }
  \draw[accent,thick,fill=white] (1,1.0) circle (1.7pt);
  \node[anchor=west,accent,font=\footnotesize] at (1.18,1.0) {members of $A=\{1-1/n\}$};
  \draw[layergrey,dashed] (1,0.95) -- (1,0);
  \node[accent,anchor=west,font=\footnotesize] at (1.12,0.45)
    {$\sup A=1$, and $1$ is not a member of $A$};
  \fill[layerteal!14] (1,-0.95) rectangle (5.35,-0.6);
  \draw[layerteal,thick] (1,-0.95) rectangle (5.35,-0.6);
  \node[layerteal!80!black,font=\footnotesize] at (3.2,-0.775) {all upper bounds of $A$};
  \draw[accent,line width=1.4pt] (1,-1.02) -- (1,-0.52);
\end{tikzpicture}
\caption{The set $A=\{1-1/n:n\in\mathbb{N}\}$ on the number line. Every $u\ge 1$
is an upper bound; the smallest, $u=1$, is $\sup A$. The number $1$ is not a member of
$A$, so $A$ has no maximum.}
\label{fig:sup-garis}
\end{figure}

\kamus{supremum}{supremum (batas atas terkecil)}{The lowest roof that still covers every
member of the set; that roof need not belong to the set.}

A supremum, if it exists, is unique; the second condition can be written with epsilon,
and that form is the one most often used in proofs.

\begin{lemma}
(i) If $A$ has a supremum, then its supremum is unique. (ii) A number $s$ is the
supremum of $A$ if and only if $s$ is an upper bound of $A$ and for every
$\varepsilon>0$ there exists $a\in A$ with $a>s-\varepsilon$.
\end{lemma}

\begin{proof}
(i) Suppose $s$ and $s'$ are both suprema. Since $s$ is an upper bound and $s'$ is a
supremum, $s'\le s$. Since $s'$ is an upper bound and $s$ is a supremum, $s\le s'$.
Hence $s=s'$.

(ii) Suppose $s$ is an upper bound and the epsilon condition holds, but there is an
upper bound $u<s$. Take $\varepsilon=s-u>0$. By the epsilon condition there is
$a\in A$ with $a>s-\varepsilon=u$, contradicting that $u$ is an upper bound.
Conversely, if $s$ is a supremum and there is $\varepsilon>0$ such that
$a\le s-\varepsilon$ for all $a\in A$, then $s-\varepsilon$ is an upper bound smaller
than $s$, contradicting that $s$ is the least upper bound.
\end{proof}

\begin{resep}
To prove that a number $s$ is the supremum of a set $A$:
\begin{enumerate}
\item Show $A\neq\varnothing$ by naming one of its members.
\item Prove that $s$ is an upper bound: take arbitrary $a\in A$, show $a\le s$.
\item Prove that no upper bound is smaller: take $\varepsilon>0$, produce
      $a\in A$ with $a>s-\varepsilon$, for instance $a=1-1/n$ with $n>1/\varepsilon$.
\item Conclude: $s$ is an upper bound and every number below $s$ is not an upper bound,
      so $s=\sup A$.
\end{enumerate}
Order warning: step 2 must be finished before step 3. Choosing $a>s-\varepsilon$ first
proves nothing, since $s$ is not yet known to be an upper bound. For
$A=\{1-1/n\}$, the number $1/2$ passes step 3 but fails step 2.
\end{resep}

The following axiom distinguishes $\mathbb{R}$ from $\mathbb{Q}$. It cannot be proved
from the field and order laws, because it is a statement about the completeness of the
line.

\begin{teorema}[Least-upper-bound property]\index{completeness axiom}
Every nonempty set $A\subseteq\mathbb{R}$ that is bounded above has a supremum in
$\mathbb{R}$.
\end{teorema}

\begin{proof}[Proof]
This statement is accepted as the completeness axiom. Other equivalent formulations:
nested interval intersections, and convergence of Cauchy sequences. We choose the
supremum because it is the most directly usable; the construction of $\mathbb{R}$ from
$\mathbb{Q}$ that proves it exists is in the Appendix to Chapter 1 of Rudin.
\end{proof}

The infimum need not be axiomatized: it is the supremum seen in a mirror.

\begin{proposisi}
If $A$ is nonempty and bounded below, then $\inf A=-\sup(-A)$, with
$-A=\{-a: a\in A\}$.
\end{proposisi}

\begin{proof}
The set $-A$ is nonempty and bounded above, so $s=\sup(-A)$ exists. For every
$a\in A$ we have $-a\le s$, hence $a\ge -s$: so $-s$ is a lower bound of $A$. Now
take $t>-s$. Then $-t<s$, so $-t$ is not an upper bound of $-A$; there is $a\in A$
with $-a>-t$, that is, $a<t$. Hence $t$ is not a lower bound of $A$. Therefore $-s$
is the greatest lower bound.
\end{proof}

First consequence of completeness: the natural numbers have no ceiling. It sounds
trivial, but without completeness it cannot be proved.

\begin{teorema}[Archimedean property]\index{archimedean property}
For every $x>0$ and every $y\in\mathbb{R}$ there exists $n\in\mathbb{N}$ with
$nx>y$.
\end{teorema}

\begin{proof}
Suppose the contrary: there are $x>0$ and $y$ such that $nx\le y$ for all $n$. Then the
set $S=\{nx: n\in\mathbb{N}\}$ is nonempty and bounded above by $y$, so
$\alpha=\sup S$ exists. Since $x>0$, the number $\alpha-x$ is smaller than $\alpha$ and
therefore not an upper bound of $S$. Hence there is $m$ with $mx>\alpha-x$, which gives
$(m+1)x>\alpha$. Yet $(m+1)x\in S$, and $\alpha$ is an upper bound of $S$. Contradiction.
\end{proof}

The reciprocal version is the one most often used: for every $\varepsilon>0$ there is
$n$ with $1/n<\varepsilon$, simply by applying the theorem to $x=\varepsilon$ and
$y=1$. The phrase ``for $n$ sufficiently large'' is shorthand for that consequence, and
now we know it is legitimate.

That property also destroys very small positive numbers: there is no $\delta>0$ that is
smaller than $1/n$ for all $n$, because $1/\delta$ would become an upper bound of
$\mathbb{N}$. Real analysis therefore needs no ``infinitesimals''.

Second consequence: every open interval, however small, contains a rational number.

\begin{teorema}[Density of $\mathbb{Q}$]\index{density}
If $x<y$ are real numbers, then there exists $p\in\mathbb{Q}$ with $x<p<y$.
\end{teorema}

\begin{proof}
Since $y-x>0$, the Archimedean property gives $n\in\mathbb{N}$ with $n(y-x)>1$. Apply
Archimedes again to obtain integers $m_1>nx$ and $m_2>-nx$; so
$-m_2<nx<m_1$. Since $nx$ is trapped between two integers, there is an integer
$m$ with $m-1\le nx<m$. Then $nx<m$, and from $m-1\le nx$ we get
$m\le nx+1<ny$. Since $n>0$, division gives $x<m/n<y$.
\end{proof}

That proof is constructive: it gives a recipe for finding $p$, namely choosing $n$
large enough and then rounding $nx$ upward. It does not give a small denominator, so
efficient approximation still demands other ideas.

\section{Why the rational numbers are not enough}
\tingkat{advanced}

The program in the opening fails because the set it squeezes has no roof in
$\mathbb{Q}$. Let us prove that, not merely display the numbers. Take
$A=\{q\in\mathbb{Q}: q^2<2\}$. This set is nonempty, since $1\in A$, and it is bounded
above inside $\mathbb{Q}$: if $q\ge 3/2$ then $q^2\ge 9/4>2$, so $3/2$ is an upper
bound.

The intervals that program keeps are drawn as a stack of ever shorter bars
(Figure~\ref{fig:selang-bersarang}). Each bar still contains root two and its length
halves, while all its endpoints stay rational; the bars narrow onto a single point
without ever touching it.

\begin{figure}[htbp]
\centering
\begin{tikzpicture}[x=9cm,y=1cm,font=\small]
  \foreach \y/\a/\b/\lab/\pan in {%
    0/1/2/{$[1,\; 2]$}/{$1$},
    0.55/1/1.5/{$[1,\; \tfrac32]$}/{$1/2$},
    1.1/1.25/1.5/{$[\tfrac54,\; \tfrac32]$}/{$1/4$},
    1.65/1.375/1.5/{$[\tfrac{11}{8},\; \tfrac32]$}/{$1/8$},
    2.2/1.375/1.4375/{$[\tfrac{11}{8},\; \tfrac{23}{16}]$}/{$1/16$},
    2.75/1.40625/1.4375/{$[\tfrac{45}{32},\; \tfrac{23}{16}]$}/{$1/32$}}{
    \fill[layerindigo!12] (\a,\y) rectangle (\b,\y+0.3);
    \draw[layerindigo,thick] (\a,\y) rectangle (\b,\y+0.3);
    \node[layerindigo,anchor=east,font=\footnotesize] at (\a-0.018,\y+0.15) {\lab};
    \node[layergrey,anchor=west,font=\footnotesize] at (\b+0.015,\y+0.15) {\pan};
  }
  \draw[accent,dashed,thick] (1.41421,-0.25) -- (1.41421,3.15);
  \node[accent,font=\footnotesize,above] at (1.41421,3.15) {$\sqrt2$};
  \draw[-{Latex[length=1.6mm]},layergrey] (1.0,-0.12) -- (2.0,-0.12);
  \node[layergrey,font=\footnotesize,below] at (1.5,-0.14) {length $=1$};
\end{tikzpicture}
\caption{The intervals kept by the opening program: $[1,2]$, $[1,\tfrac32]$,
$[\tfrac54,\tfrac32]$, $[\tfrac{11}{8},\tfrac32]$, $[\tfrac{11}{8},\tfrac{23}{16}]$,
$[\tfrac{45}{32},\tfrac{23}{16}]$. Their lengths fall from $1$ to $1/32$; continued
without end, the intersection shrinks to a single point, namely $\sqrt2$, which is not
a rational number.}
\label{fig:selang-bersarang}
\end{figure}

First step: there is no largest member. Every rational with $q^2<2$ can still be nudged
a little upward without passing two, and the size of the nudge has a formula.

\begin{lemma}
If $q\in\mathbb{Q}$, $q>0$, and $q^2<2$, then there exists $h\in\mathbb{Q}$, $h>0$,
with $(q+h)^2<2$. In particular $q+h\in A$ and $q+h>q$.
\end{lemma}

\begin{proof}
Take $h=\dfrac{2-q^2}{2q+2}$. Since $q^2<2$ and $2q+2>0$, we have $h>0$ and $h$
rational. Also $h<1$, because $2-q^2<2q+2$ is equivalent to $q^2+2q>0$, which is true
for $q>0$. Then $h^2\le h$, so
\[
(q+h)^2=q^2+2qh+h^2\le q^2+h(2q+1)
=q^2+(2-q^2)\frac{2q+1}{2q+2}<q^2+(2-q^2)=2,
\]
because $\frac{2q+1}{2q+2}<1$.
\end{proof}

That formula can be checked numerically. For $q=1393/985$ we have $2-q^2=1/970225$, so
$h=1/4684660$ and
\[
q+h=\frac{6625109}{4684660}=1.414213411432206\ldots,
\]
whose square is $1.999999573074719\ldots<2$. So $1393/985$ is not the largest member,
even though its squared error is already about one part in a million. This procedure
can be repeated without end: there is no last step.

\begin{tangan}
Materials: 1.5 m of rope, six pieces of tape, a tape measure.

Stick the tape on the floor perpendicular to the wall, at distances $60/n$ cm for
$n=1,\ldots,6$: 60, 30, 20, 15, 12, and 10 cm. You have only stuck six, but the rule
does not stop. Stretch the rope across as a fence, then slide it toward the wall. Each
time you stop, check one thing: are all the tape marks on the same side of the rope as
the wall? As long as they are, the rope is still a fence.

The rope can be brought arbitrarily close to the wall and still be a fence: 5 cm, 1 cm,
1 mm. What stops it is not the tape but the wall. The collection of tape marks has no
smallest member, and not one of them touches the wall. The lowest fence still exists
--- the wall --- and it is not a member of the collection. That is the supremum.
\end{tangan}

It is not only the maximum that is missing. The least upper bound is missing from the
rational numbers as well, and the proof requires one fact about root two.

\begin{proposisi}
The set $A=\{q\in\mathbb{Q}: q^2<2\}$ has no supremum in $\mathbb{Q}$.
\end{proposisi}

\begin{proof}
Suppose $p\in\mathbb{Q}$ is the supremum of $A$. Since $1\in A$, we have $p\ge 1>0$. We
show that $p^2=2$, which is impossible because no rational number has square two.

If $p^2<2$, the previous lemma gives $q=p+h\in A$ with $q>p$, contradicting that $p$
is an upper bound. If $p^2>2$, take $k=(p^2-2)/(2p)>0$, which is rational and smaller
than $p$ since $p^2-2<2p^2$. For every $q\in A$,
\[
(p-k)^2=p^2-2pk+k^2>p^2-2pk=2>q^2,
\]
so $q<p-k$ for $q>0$, and for $q\le 0$ this is clear because $p-k>0$: the number
$p-k$ is still an upper bound of $A$, but smaller than $p$.
That contradicts $p$ being the least upper bound. Hence $p^2=2$.

Now $p=m/n$ with $m,n$ integers having no common factor. From $m^2=2n^2$, $m^2$ is
even, so $m$ is even; write $m=2r$. Then $4r^2=2n^2$, that is, $n^2=2r^2$, so $n$
is even as well. Both are even, contradicting the choice of $m,n$.
\end{proof}

Inside $\mathbb{R}$ that does not happen. The same set, viewed this time as a set of
real numbers, is bounded above by $2$, so by the completeness axiom it has a supremum.
That supremum is the number we name $\sqrt{2}$. The following proof closes the two
remaining possibilities, and its shape imitates the earlier lemma with $q$ replaced by
a real number.

\begin{teorema}[Existence of root two]\index{square root of two}
There is exactly one $s>0$ with $s^2=2$.
\end{teorema}

\begin{proof}
Let $E=\{t\in\mathbb{R}: t>0,\ t^2<2\}$. This set is nonempty because $1\in E$, and it
is bounded above by $2$: if $t>2$ then $t^2>4>2$, so $t\notin E$. By the
least-upper-bound property, $s=\sup E$ exists, and $s\ge 1>0$ because $1\in E$.

Suppose $s^2<2$. Take $h=\frac12\min\{1,(2-s^2)/(2s+1)\}$. Then $0<h<1$ and
$h(2s+1)<2-s^2$: if the first value is the minimum, $h=\frac12$ and
$\frac12(2s+1)\le\frac12(2-s^2)<2-s^2$; if the second is, $h(2s+1)=\frac12(2-s^2)$.
Since $h<1$, we have $h^2<h$, so
\[
(s+h)^2=s^2+2sh+h^2<s^2+h(2s+1)<s^2+(2-s^2)=2 .
\]
Hence $s+h\in E$ although $s+h>s$, contradicting that $s$ is an upper bound of $E$.
Therefore $s^2\ge 2$.

Suppose $s^2>2$. Take $k=(s^2-2)/(2s)>0$; then $k<s$ because $s^2-2<2s^2$. For every
$t\in E$ with $t\ge s-k$ we have
\[
t^2\ge(s-k)^2=s^2-2sk+k^2>s^2-2sk=2,
\]
so $t\notin E$. Hence $s-k$ is an upper bound of $E$ smaller than $s$, contradicting
that $s$ is the least upper bound. Since $s^2<2$ and $s^2>2$ are both impossible,
$s^2=2$.

Uniqueness: if $s,t>0$ and $s^2=t^2=2$, then $(s-t)(s+t)=0$. Since $s+t>0$, $s=t$.
\end{proof}

The number $s$ is irrational, by the parity argument above: the hole in $\mathbb{Q}$ is
real, and $\mathbb{R}$ closes it with exactly one number. What closes it is not order
and not the field laws, but completeness. The opening program fails not from lack of
precision, but because its number system does not contain the object it is looking for.

One thing must be stated honestly: we have not built $\mathbb{R}$ from $\mathbb{Q}$, we
have only assumed that $\mathbb{R}$ satisfies three groups of axioms. The construction
is known. Dedekind cuts $\mathbb{Q}$ in two: a \emph{cut} is a subset
$\alpha\subset\mathbb{Q}$ that is nonempty, not all of $\mathbb{Q}$, closed downward,
and without a largest member. That last property is exactly our lemma above. Cantor
uses equivalence classes of rational Cauchy sequences; both appeared in 1872. The
details are in the Appendix to Chapter 1 of Rudin.

That construction explains the standing of decimals. For $x>0$, take $n_0$ the largest
integer with $n_0\le x$, then $n_1$ the largest digit with $n_0+n_1/10\le x$, then
$n_2$ with $n_0+n_1/10+n_2/100\le x$, and so on. That cut is bounded above by $x$, so
its supremum exists and equals $x$. A decimal is not the raw material of the real
numbers but a way of naming a supremum; completeness guarantees that the name exists,
not that the name can be written with fractions.

\kamus{completeness axiom}{aksioma kelengkapan}{The promise that the number line has no
holes: every collection that is bounded above always has a lowest roof.}

The existence of $\sqrt2$ makes the computational problem sensible: $s$ can be
computed without knowing its digits in advance. Newton's iteration,
$x_{k+1}=\frac12(x_k+2/x_k)$, stays rational when $x_k$ is rational, so the whole
process runs inside $\mathbb{Q}$. From $s^2=2$ we get
\[
x_{k+1}-s=\frac{(x_k-s)^2}{2x_k},\qquad e_{k+1}=e_k^2\cdot\frac{s}{2x_k},
\]
with $e_k=|x_k-s|/s$. The right side of the first form is nonnegative, so every step
stays above $s$; there the factor $\frac{s}{2x_k}$ does not exceed $1/2$, so
$e_{k+1}\le e_k^2/2$ and the number of correct digits roughly doubles at each step.

One warning: we have not proved that $x_k\to s$. That needs the tools of Chapter 4.
What can be said now: since $e_1<0.061$, from $e_{k+1}\le e_k^2/2$ one can derive
$e_k\le 2^{-(2^k-1)}$ for $k\ge 1$, and that is enough to design an honest numerical
test.

\section{Intervals, accumulation points, and open sets}
\tingkat{advanced}

With the supremum in hand, we can name the pieces of the number line, like the finish
rules of a race: the line is the same, only whether the line itself counts differs.
Those names are used again in Chapter 11 and Chapter 13, so we take only what we need
without building a topology in full.

An interval is a piece of the line with no holes. The following definition writes the
nine possible forms; all are determined by two endpoints and by the decision whether
each endpoint counts.

\begin{definisi}[Interval]
Let $a<b$ be real numbers. The set $(a,b)=\{x: a<x<b\}$ is called an open interval,
$[a,b]=\{x: a\le x\le b\}$ a closed interval, and $[a,b)$, $(a,b]$ are mixed forms. The
sets $(a,\infty)=\{x: x>a\}$, $[a,\infty)$, $(-\infty,b)$,
$(-\infty,b]$, and $\mathbb{R}$ itself are called unbounded intervals. All nine forms
are called intervals.
\end{definisi}

The feature that marks them all: if two points lie in the interval, the whole segment
between them lies there too, so intervals are called hole-free. The list above is not
proved complete; the classification only inspects upper bound, lower bound, and whether
that bound is included. The notion of a point being ``near'' is expressed by a
neighborhood, that is, a small open interval around one point.

\begin{definisi}[Neighborhood, interior point, open, closed]\index{open set}\index{closed set}\index{accumulation point}
For $\varepsilon>0$, the neighborhood of a point $x$ is
$V_\varepsilon(x)=(x-\varepsilon,x+\varepsilon)$. A point $x$ is called an interior
point of $S\subseteq\mathbb{R}$ if $V_\varepsilon(x)\subseteq S$ for some
$\varepsilon>0$; $S$ is called open if every member is an interior point. A point $x$ is
called an accumulation point of $S$ if every $V_\varepsilon(x)$ contains a member of
$S$ different from $x$. A set $S$ is called closed if its complement is open.
\end{definisi}

A few examples organize the intuition. The interval $(0,1)$ is open and not closed,
because every neighborhood of $0$ meets it. The interval $[0,1]$ is closed and not open,
while $[0,1)$ is neither. The set $\mathbb{Q}$ is also neither, because every
neighborhood contains rationals and irrationals. The set $\{1/n:n\in\mathbb{N}\}$ is not
closed: $0$ lies in its complement, but every neighborhood of $0$ contains some $1/n$.
The set $\mathbb{Z}$ is closed, because its complement is the union of the open
intervals $(k,k+1)$.

So ``closed'' is not the negation of ``open'': $\mathbb{R}$ and $\varnothing$ are both
open and closed at once, and the complement of an open set is always closed. The density
of $\mathbb{Q}$ makes these examples sharper. Every neighborhood contains rationals and
irrationals, so no point is an interior point of $\mathbb{Q}$ or of its complement: that
set is dense and hole-ridden everywhere at once. Completeness therefore cannot be
replaced by density.

A neighborhood is itself open, and that is what makes the definition of an interior
point reusable.

\begin{proposisi}
(i) Every neighborhood $V_\varepsilon(x)$ is open. (ii) The union of an arbitrary
family of open sets is open. (iii) The intersection of finitely many open sets is open.
\end{proposisi}

\begin{proof}
(i) Take $y\in V_\varepsilon(x)$, so $|y-x|<\varepsilon$. Set
$\delta=\varepsilon-|y-x|>0$. For every $z\in V_\delta(y)$,
$|z-x|\le|z-y|+|y-x|<\delta+|y-x|=\varepsilon$ by the triangle inequality. Hence
$V_\delta(y)\subseteq V_\varepsilon(x)$.

(ii) If $x$ lies in the union, it lies in one of its members, and since that member is
open there is a neighborhood of $x$ contained in it, hence also contained in the union.

(iii) If $x$ lies in the finite intersection $S_1\cap\cdots\cap S_n$, then for each $i$
there is $\varepsilon_i>0$ with $V_{\varepsilon_i}(x)\subseteq S_i$. Take
$\varepsilon=\min\{\varepsilon_1,\ldots,\varepsilon_n\}>0$. Then $V_\varepsilon(x)$ is
contained in every $S_i$, hence contained in the intersection.
\end{proof}

Item (iii) holds only for finite intersections: $\bigcap_n(-1/n,1/n)=\{0\}$ is not open,
because for $x\neq0$ the Archimedean property gives $n$ with $1/n<|x|$.

\section{Worked examples}
\tingkat{intermediate}

Three computations promised earlier are carried out here: Newton's iteration, rational
approximation to error below $10^{-6}$, and a check that $\{q:q^2<2\}$ has no maximum.
All the numbers are computed with exact fractions.

Starting from $x_0=1$, the iteration $x_{k+1}=(x_k+2/x_k)/2$ gives the following
fractions; the denominators grow, and that is the price paid for accuracy.

\begin{center}\small
\begin{tabular}{rlll}
\hline
$k$ & $x_k$ & decimal & relative error \\
\hline
0 & $1$ & $1$ & $2.92893\times10^{-1}$ \\
1 & $3/2$ & $1.5$ & $6.06602\times10^{-2}$ \\
2 & $17/12$ & $1.41666666667$ & $1.73461\times10^{-3}$ \\
3 & $577/408$ & $1.41421568627$ & $1.50183\times10^{-6}$ \\
4 & $665857/470832$ & $1.41421356237$ & $1.12774\times10^{-12}$ \\
5 & $886731088897/627013566048$ & $1.41421356237$ & $6.35896\times10^{-25}$ \\
\hline
\end{tabular}
\end{center}

The relative error is measured against $\sqrt{2}=1.414213562373095\ldots$. The reduction
factor climbs by itself: from step 0 to 1 about $4.8$, then $35$, $1155$, and
$1.3\times10^{6}$ --- that is the signature of squaring. The ratio $e_{k+1}/e_k^2$ takes
the values $0.7071$, $0.4714$, $0.4991$, $0.49999925$, and $0.5$, in agreement with the
formula $e_{k+1}=e_k^2\cdot\frac{s}{2x_k}$ because $x_k$ approaches $s$ from above. At
steps 4 and 5 the decimals agree to eleven digits; the next digits do not appear because
we print twelve significant digits, not because the error stops falling. The fraction at
step 5 has twelve digits with error $6.36\times10^{-25}$, and one additional digit
demands roughly one more step.

The second method uses rational approximations that spend denominators more
economically, namely the continued fraction of $\sqrt2$:
\[
1,\quad \frac32,\quad \frac75,\quad \frac{17}{12},\quad \frac{41}{29},\quad \ldots
\]
with $p_{k+1}=p_k+2q_k$ and $q_{k+1}=p_k+q_k$ starting from $(p_0,q_0)=(1,1)$. The
following table lists their absolute errors.

\begin{center}\small
\begin{tabular}{rlll}
\hline
$p/q$ & decimal & absolute error & $p^2-2q^2$ \\
\hline
$1/1$ & $1$ & $4.14214\times10^{-1}$ & $-1$ \\
$3/2$ & $1.5$ & $8.57864\times10^{-2}$ & $1$ \\
$7/5$ & $1.4$ & $1.42136\times10^{-2}$ & $-1$ \\
$17/12$ & $1.41666666667$ & $2.45310\times10^{-3}$ & $1$ \\
$41/29$ & $1.41379310345$ & $4.20459\times10^{-4}$ & $-1$ \\
$99/70$ & $1.41428571429$ & $7.21519\times10^{-5}$ & $1$ \\
$239/169$ & $1.41420118343$ & $1.23789\times10^{-5}$ & $-1$ \\
$577/408$ & $1.41421568627$ & $2.12390\times10^{-6}$ & $1$ \\
$1393/985$ & $1.41421319797$ & $3.64404\times10^{-7}$ & $-1$ \\
$3363/2378$ & $1.41421362489$ & $6.25218\times10^{-8}$ & $1$ \\
\hline
\end{tabular}
\end{center}

The first approximation whose error is below $10^{-6}$ is $1393/985$ with error
$3.64\times10^{-7}$; the one before it, $577/408$, is still $2.12\times10^{-6}$. The last
column explains why: since $p^2-2q^2=\pm1$,
\[
\left|\frac pq-\sqrt2\right|
=\frac{|p^2-2q^2|}{q^2\left(\frac pq+\sqrt2\right)}
<\frac{1}{2\sqrt2\,q^2},
\]
so the error falls like the inverse square of the denominator. For $q=985$ that bound is
$3.644\times10^{-7}$ and the true error is $3.64404\times10^{-7}$: almost exactly,
because $\frac pq$ is already very close to $\sqrt2$. The denominator 985 is far smaller
than $627013566048$ at Newton step 5, although its error is only about a million times
larger --- that is the difference between a method chasing quadratic error and one
chasing the smallest rational error.

Figure~\ref{fig:tangga-akar} displays that error column as a descending staircase. The
sign of $p^2-2q^2$ alternates, so the fractions jump from below to above root two, and
the ratio of two successive errors approaches $5.8284=(1+\sqrt2)^2$.

\begin{figure}[htbp]
\centering
\begin{tikzpicture}[x=1.02cm,y=1cm,font=\small]
  \foreach \lv in {1,...,7}{
    \draw[layergrey!20] (0.55,\lv) -- (10.75,\lv);
    \node[layergrey,anchor=east,font=\scriptsize] at (0.45,\lv) {$10^{-\lv}$};
  }
  \node[layergrey,anchor=east,font=\scriptsize] at (0.45,0) {$10^{0}$};
  \draw[-{Latex[length=2mm]},ink,thick] (0.55,0) -- (10.95,0);
  \draw[-{Latex[length=2mm]},ink,thick] (0.55,0) -- (0.55,7.85);
  \foreach \k/\lab in {%
    1/{$1/1$},2/{$3/2$},3/{$7/5$},4/{$17/12$},5/{$41/29$},
    6/{$99/70$},7/{$239/169$},8/{$577/408$},9/{$1393/985$},10/{$3363/2378$}}{
    \draw[ink] (\k,0.07) -- (\k,-0.07);
    \node[rotate=32,anchor=north east,font=\scriptsize,layergrey] at (\k,-0.12) {\lab};
  }
  \foreach \k/\v in {1/0.3828,3/1.8473,5/3.3763,7/4.9073,9/6.4384}{
    \fill[accent] (\k,\v) circle (1.8pt);
  }
  \foreach \k/\v in {2/1.0666,4/2.6103,6/4.1418,8/5.6729,10/7.2040}{
    \fill[layerteal] (\k,\v) circle (1.8pt);
  }
  \draw[layergrey,dashed,thick] (0.55,6) -- (10.75,6);
  \node[layergrey,font=\scriptsize,anchor=south east] at (10.7,6.02) {threshold $10^{-6}$};
  \fill[accent] (0.95,7.55) circle (1.8pt);
  \node[right,font=\scriptsize] at (1.05,7.55) {below $\sqrt2$};
  \fill[layerteal] (4.35,7.55) circle (1.8pt);
  \node[right,font=\scriptsize] at (4.45,7.55) {above $\sqrt2$};
  \node[font=\scriptsize,anchor=north] at (5.7,-1.18) {successive fractions};
  \node[font=\scriptsize,rotate=90,anchor=south] at (-0.35,3.9) {$|p/q-\sqrt2|$};
\end{tikzpicture}
\caption{A staircase of rational approximations: the error $|p/q-\sqrt2|$ for ten
successive fractions, logarithmic scale. The $10^{-6}$ threshold is first crossed by
$1393/985$, with error $3.644\times10^{-7}$; after that each step falls by a factor of
$5.8284$.}
\label{fig:tangga-akar}
\end{figure}

The third check concerns the absence of a maximum. One more step from the previous $q$
gives
\[
q'=\frac{149856914148961}{105964844887080}=1.41421349985133266\ldots,
\]
and its square is still below two. Both numbers are rational, both lie in $A$, and each
is larger than the one before: there is no reason for this process to stop.

\begin{tantangan}
Consider $S=\{q\in\mathbb{Q}: q\ge 0,\ q^2\le 2\}$. At first glance this set has a
largest member: the number whose square is exactly two, since the sign $\le$ seems to
include that bound in the set. Prove that $S$ has no maximum, then determine $\sup S$ in
$\mathbb{R}$ and explain why the answer is not a member of $S$.

Hint 1: members of $S$ must be rational. First check whether your candidate for largest
is even allowed to be a member of the set.

Hint 2: for every $q\in S$ with $q>0$ we have $q^2<2$ strictly, because $q^2=2$ is
impossible for a rational. So the shift lemma $h=(2-q^2)/(2q+2)$ applies and produces a
larger member of $S$; the unattainable bound is $\sqrt2$.

If you get stuck, go back to the section ``Why the rational numbers are not enough'' and
to the lemma on shifting.
\end{tantangan}

Finally, the Archimedean property in numbers. For $\varepsilon=10^{-6}$ we need
$n>10^{6}$. At $n=10^{6}$ the reciprocal is exactly $10^{-6}$, which is not less than
$\varepsilon$; at $n=10^{6}+1$ the reciprocal is $9.99999\times10^{-7}$, which does
qualify. One unit in $n$ decides whether the choice is legitimate, and that is why
strict inequalities must be respected.

\kertas{For $x=1393/985$, compute $h=(2-x^2)/(2x+2)$ and check by hand that
$(x+h)^2<2$. Then compute the relative error $|x-\sqrt2|/\sqrt2$ and compare it with
$1/(2\sqrt2\,x^2)$. Explain why the two are almost equal, and why that fact does not
contradict the absence of a maximum in $A$.}

\section{Computational lab}
\tingkat{advanced}

The following code runs with Python 3 and \texttt{mpmath} for high-precision
arithmetic; Newton's iteration uses \texttt{Fraction}, so there is no rounding before
printing.

\begin{lab}
\begin{verbatim}
from fractions import Fraction as F
import mpmath as mp
mp.mp.dps = 40
root2 = mp.sqrt(2)

print("Newton for root two, x0 = 1")
x = F(1)
for k in range(6):
    v = mp.mpf(x.numerator) / x.denominator
    print(k, f"{x.numerator}/{x.denominator}", mp.nstr(v, 12),
          mp.nstr(abs(v - root2) / root2, 6))
    x = (x + F(2) / x) / 2

print()
print("Successive rational approximations, error < 1e-6 ?")
p, q = 1, 1
for k in range(10):
    v = mp.mpf(p) / q
    print(f"{p}/{q}", mp.nstr(v, 12), mp.nstr(abs(v - root2), 6),
          abs(v - root2) < mp.mpf(10) ** -6)
    p, q = p + 2 * q, p + q

print()
print("Test: {rational q : q^2 < 2} has no maximum")
q = F(1393, 985)
for k in range(2):
    h = (2 - q * q) / (2 * q + 2)
    q = q + h
    print(q, mp.nstr(mp.mpf(q.numerator) / q.denominator, 18), q * q < 2)
\end{verbatim}
\end{lab}

The output obtained, verbatim:

\begin{verbatim}
Newton for root two, x0 = 1
0 1/1 1.0 0.292893
1 3/2 1.5 0.0606602
2 17/12 1.41666666667 0.00173461
3 577/408 1.41421568627 1.50183e-6
4 665857/470832 1.41421356237 1.12774e-12
5 886731088897/627013566048 1.41421356237 6.35896e-25

Successive rational approximations, error < 1e-6 ?
1/1 1.0 0.414214 False
3/2 1.5 0.0857864 False
7/5 1.4 0.0142136 False
17/12 1.41666666667 0.0024531 False
41/29 1.41379310345 0.000420459 False
99/70 1.41428571429 7.21519e-5 False
239/169 1.41420118343 1.23789e-5 False
577/408 1.41421568627 2.1239e-6 False
1393/985 1.41421319797 3.64404e-7 True
3363/2378 1.41421362489 6.25218e-8 True

Test: {rational q : q^2 < 2} has no maximum
6625109/4684660 1.41421341143220639 True
149856914148961/105964844887080 1.41421349985133266 True
\end{verbatim}

Three things are readable from that output. The Newton error stops being printed at
twelve digits at step 4, even though the true value still falls to $6.4\times10^{-25}$:
what limits it is the print width, not the mathematics. The $10^{-6}$ threshold is
crossed exactly at $1393/985$, where the column of \texttt{True} begins. Both numbers in
the maximum test are also \texttt{True}: both are rational and both have squares below
two.

That code proves nothing. It checks that the formulas in this chapter give the numbers
stated and that the direction of the inequalities is right; the proofs remain in the
lemmas and theorems. Program output is data, not argument.

\begin{untukdosen}
The teaching order I suggest: the root-finding program in the opening, then the field
axioms with $0\cdot x=0$ proved on the board, then order, absolute value, and
completeness. Do not postpone completeness; the entire content of the chapter depends on
it.

Two misconceptions appear often. First, students think a supremum must be a member of
the set, so they reject $\sup\{1-1/n\}=1$. Second, students think completeness can be
proved from the field and order laws, whereas $\mathbb{Q}$ satisfies both and fails
completeness. Both disappear once students see those two examples side by side.

A blackboard activity that works: each student takes $q=1.4$, computes
$h=(2-q^2)/(2q+2)$ and checks $(q+h)^2<2$; the differing results are precisely what
shows the formula works for every $q$. A readiness signal: students can explain why $A$
has no maximum without computing.
\end{untukdosen}

\section{Summary}
\tingkat{basic}

\begin{ringkasan}
\begin{itemize}
\item The real numbers rest on three groups of axioms: field, order, completeness.
      The first gives algebra, the second comparison, the third closes the holes.
\item From the field laws we derive $0\cdot x=0$, $(-x)y=-(xy)$, and the absence of
      zero divisors.
\item From the order laws we derive the sign rules, including nonnegative squares and
      $1>0$.
\item Absolute value turns distance into inequality: $|x|\le c$ is equivalent to
      $-c\le x\le c$, and $|x+y|\le|x|+|y|$.
\item Completeness guarantees a supremum for every nonempty set that is bounded above,
      and then yields the Archimedean property and the density of $\mathbb{Q}$.
\item $\mathbb{Q}$ satisfies the field and order laws but is not complete:
      $\{q: q^2<2\}$ has no maximum; in $\mathbb{R}$ its supremum is $\sqrt{2}$.
\end{itemize}
\end{ringkasan}

\section{Exercises}
\tingkat{intermediate}

\begin{latihan}
\begin{enumerate}
  \item Prove that in an ordered field, $a<b$ is equivalent to $-b<-a$, and conclude
        that multiplying an inequality by $-1$ reverses its direction.
  \item Prove $a<\frac{a+b}{2}<b$ for $a<b$, then explain why that result does not
        require completeness whereas $\sup\{q:q^2<2\}$ does.
  \item Compute $\sup$ and $\inf$ of the set $S=\{2-1/n : n\in\mathbb{N}\}$, then check
        both conditions of the definition. Also determine whether $S$ has a maximum.
  \item Assess this claim: ``Since every positive number is greater than $1/n$ for some
        $n$, there is no smallest positive real number.'' True, false, or not yet
        proved? Name the part that the proof supports.
  \item With $h=(2-q^2)/(2q+2)$, compute $q+h$ for $q=1$ and $q=7/5$, then check that
        both still lie in $A=\{q\in\mathbb{Q}:q^2<2\}$. Compare the sizes of the
        increases.
  \item Show that if $|x|\le 2$ and $|y|\le 2$ then
        $|x^3-y^3|\le 12\,|x-y|$, and say where the triangle inequality is used.
  \item Prove that if $S$ is open then its complement is closed, and conversely, using
        the definition of a neighborhood.
  \item Compute five steps of Newton's iteration for $\sqrt{3}$ starting from $x_0=2$,
        and report the relative error at each step. Check whether the ratio
        $e_{k+1}/e_k^2$ approaches $1/2$.
  \item Show $\bigcap_{n=1}^{\infty}(-1/n,1/n)=\{0\}$ and explain why that does not
        contradict the closedness of finite intersections.
  \item Prove that between any two real numbers there is always an irrational number.
\end{enumerate}

\noindent\textbf{Hints.} Item 4: the claim contains two different statements, one about
every $\varepsilon$ and one about the existence of a minimum. Item 8: use
$e_{k+1}=e_k^2\cdot\frac{\sqrt3}{2x_k}$ and $x_k\ge\sqrt3$. Item 10: apply the density
of $\mathbb{Q}$ to $x/\sqrt2$ and $y/\sqrt2$, then multiply back.
\end{latihan}

\begin{untukpeneliti}
What is genuinely proved: the completeness axiom does not follow from the field and
order laws, because $\mathbb{Q}$ satisfies both and fails the third. The existence of a
complete ordered field is proved by Rudin via Dedekind cuts, and there it is also stated
without proof that any two complete ordered fields are isomorphic; that claim is not
used here.

Primary sources: Rudin, \emph{Principles of Mathematical Analysis}, Chapter 1
(Definition 1.12, Proposition 1.18, Theorems 1.19--1.21, the Appendix on cuts); Abbott,
\emph{Understanding Analysis}, Chapter 1 (The Axiom of Completeness and its
consequences).

What remains open: the effective content of Newton's iteration. Without an explicit
initial bound, $e_{k+1}\le e_k^2/2$ does not yet give a step count; that requires
Chapter 4 and, for general functions, Kantorovich theory. The chapter's scope: topology
only as far as intervals and neighborhoods, and $\mathbb{R}$ assumed rather than built.
\end{untukpeneliti}

\peta{Chapter 3 in Rudin, \emph{Principles of Mathematical Analysis}: Chapter 1 (Ordered
Sets, Fields, The Real Field, The Least-Upper-Bound Property, the Archimedean property,
density of $\mathbb{Q}$, existence of $n$-th roots, Appendix on Dedekind cuts). In
Abbott, \emph{Understanding Analysis}: Chapter 1 (The Axiom of Completeness and
Consequences of Completeness); open sets and accumulation points from Chapter 3, not
developed here.}
